\paragraph{The semi-classical criterion and the signature of the kinetic tensor on the supplied backgrounds (conditional; bounded no-go).}
The preceding paragraph concerns the contraction \(h_0\) of the matrix indices, not the kinetic tensor of a background. This paragraph treats the latter for the first-order equivariant class of Section~\ref{sec:status-D} and the conjugate-coordinate class of Section~\ref{sec:status-F}, and tests both against the semi-classical criterion of~\cite{Steinacker2026Weak}. Keep Equation~\eqref{eq:Smatrix} with \(C_{\mu\nu\lambda}=0\), a real contraction \(h_0\) and, where stated, the mass term \(S_\gamma\) contracted with \(h_0\), whose equation of motion is \([A^\mu,[A_\mu,A_\nu]]=\gamma A_\nu\).

\emph{Criterion.} For a background \(\bar T^a\), reference~\cite{Steinacker2026Weak} takes quasi-coherent states \(|t\rangle\), the ground states of \(\sum_a(\bar T^a-t^a)^2\), and their uncertainty \(\Delta_T=\max_a\langle t|(\bar T^a-\mathbf t^a)^2|t\rangle^{1/2}\) with \(\mathbf t^a=\langle t|\bar T^a|t\rangle\), and proposes as the condition for the semi-classical regime that \(\max_{a,t'}\bigl|\bigl\langle\langle t|\mathcal A_a|t'\rangle\langle t'|\mathcal A_a|t\rangle\bigr\rangle\bigr|^{1/2}\ll\Delta_T\) for the fluctuations \(\mathcal A\) under the matrix integral, with \(|t'-t|<\Delta_T\) for Minkowski signature (its Eqs.~(13), (16) and~(17)). It is evaluated there with the tree-level propagator in a Lorentz gauge, and it is stated there to fail for commutative backgrounds.

\emph{The flat point of the first-order equivariant class.} There \(\hat X_a\) is multiplication by \(-\tilde p_a\), \(\tilde p=up\). The operator \(\sum_a(\hat X_a-t_a)^2\) has no eigenvector, and normalised states concentrated at \(\tilde p=-t\) have \(\Delta_T\to0\); for commuting \(N\times N\) matrices the eigenvectors give \(\Delta_T=0\) exactly, while the off-diagonal entries fluctuate. The criterion therefore fails, as stated in~\cite{Steinacker2026Weak} for commutative backgrounds. For a normalised state whose density in \(\tilde p\) is Gaussian with variance \(\sigma^2\) in each component, so that \(\Delta_T=\sigma\), the background-gauge form of that section gives at \(t'=t\), exactly and with the fibre integrated out by the unit Jacobian of \(p\mapsto up\), \(\langle|\langle t|\mathcal A_a|t\rangle|^2\rangle=g_{\mathrm{YM}}^2/4\sigma^2\) for the Euclidean contraction, and a value of modulus \(g_{\mathrm{YM}}^2/(2\sqrt2\,\sigma^2)\) for \(\eta\) with Definition~(A). The ratio of the two sides of the criterion is then \(g_{\mathrm{YM}}/2\sigma^2\), respectively \(2^{-3/4}g_{\mathrm{YM}}/\sigma^2\); for the Euclidean contraction the maximum is at \(t'=t\), and for \(\eta\) no larger value occurs at the sampled \(|t'-t|\le\sigma\) (numerical). So a semi-classical regime requires states of width \(\sigma\gg g_{\mathrm{YM}}^{1/2}\), a scale which the commuting background does not fix. The quantity \(g_{\mathrm{YM}}/\sigma^2\) has the form of the coupling \(g/\Delta^2\) of the Moyal plane in~\cite{Steinacker2026Weak}, with the width of the state in place of the noncommutativity scale; the weight \(w\) used for this class in Section~\ref{sec:status-D} supplies such a width and is not derived.

\emph{Signature of the kinetic tensor.} For a frame \(E^{a\mu}\), \(\{\bar T^a,y^\mu\}\) in the Poisson limit or \(\mathbb D_ay^\mu\) for first-order operators, the kinetic tensor is \(\gamma^{\mu\nu}=(h_0)_{ab}E^{a\mu}E^{b\nu}\). (i) If \(h_0\) is positive semidefinite on the span of the \(E^{\cdot\mu}\), then \(\gamma\) is a Gram matrix and positive semidefinite, and so is every average of such tensors with a positive weight; for the Euclidean contraction no kinetic tensor of this form is Lorentzian. (ii) Where \(E\) has rank four and does not depend on the fibre point, \(\gamma\) has, by Sylvester's law of inertia, the signature of the restriction of \(h_0\) to the four-plane \(W\) spanned by the \(E^{\cdot\mu}\); for \(h_0\) of signature \((9,1)\) it is Lorentzian, positive definite or degenerate according as \(W\) contains a timelike vector, is spacelike or is null. (iii) An average over the fibre of tensors of signature \((3,1)\) can have another signature.

\emph{Conjugate coordinates.} With \(T^a=M^{a4}\), \(y^\mu=M^{\mu5}\), \(m=M^{45}\) and \(\ell^2=-y_ay^a\), one has \(E^{a\mu}=\pm m\,\eta^{a\mu}\) at every point, independent of the fibre (Section~\ref{sec:status-F}). If the background is \(\bar T^I=L^I{}_aT^a\) with a real \(10\times4\) matrix \(L\) of rank four, then \(\gamma=m^2\eta^{-1}H\eta^{-1}\) with \(H=L^{\mathrm T}h_0L\): by (ii) its signature is that of \(H\) wherever \(m\neq0\), and \(\gamma=0\) where \(m=0\), the set on which the frame vanishes. For generators \(T^a=M^{ac_0}\) of \(\mathfrak{so}(p,q)\) with diagonal invariant form \(\eta\) and \(a\) in a set of four indices other than \(c_0\), the Lie algebra gives exactly
\[
H_{cc'}\,[T^c,[T^{c'},T^a]]=\eta^{c_0c_0}\bigl[\operatorname{tr}(\eta^{-1}H)\,T^a-(\eta^{-1}H)^a{}_b\,T^b\bigr]
\]
(symbolic check for every symmetric \(H\), in \(\mathfrak{so}(4,2)\), \(\mathfrak{so}(4,1)\) and \(\mathfrak{so}(6)\)). So \(\bar T^I\) solves the equation of motion exactly when \(H=c\,\eta\) on the four indices, with \(\gamma=3c\,\eta^{c_0c_0}\). For the class above, \(c_0=4\) and \(\eta=\mathrm{diag}(-1,1,1,1)\) on the four indices, so for \(h_0\) of signature \((9,1)\) this forces \(c>0\): \(W\) is timelike, \(\gamma>0\), and the kinetic tensor is Lorentzian on the whole set \(m\neq0\), and the \(h_0\)-orthogonal complement of \(W\), in which the factor \(K\) of the product backgrounds is placed, is positive definite. The same identity for \(\mathfrak{so}(6)\), with \(\eta=\delta\) and the same formula for \(E\), which holds on every coadjoint orbit, gives at the same sign \(\gamma>0\) a positive definite kinetic tensor on \(m\neq0\) for spacelike \(W\). A representation of \(\mathfrak{so}(4,2)\) by Hermitian \(N\times N\) matrices is trivial, since its image would be a compact simple Lie algebra isomorphic to \(\mathfrak{so}(4,2)\), whose Killing form is indefinite, whereas \(\mathfrak{so}(6)\) has such representations. Within this family, the non-zero solutions at finite \(N\) are therefore those of compact forms, with a positive definite kinetic tensor, and the Lorentzian member exists only for \(N\to\infty\).

\emph{Base coordinates as matrices.} If instead the four matrices are the base functions, \(\bar T^\mu=y^\mu\), the identity with \(c_0=5\) gives \(\gamma=-3c\), and \(\{y^\mu,y^\nu\}=\pm M^{\mu\nu}\) depends on the fibre point. On the orbit \(\eta_{\alpha\beta}M^{\alpha\mu}M^{\beta\nu}=\ell^2\eta^{\mu\nu}+y^\mu y^\nu-t^\mu t^\nu\), \(t^\mu=T^\mu\), which at a fibre point has signature \((3,1)\) with its negative direction along \(t\), a spacelike vector. Averaged over the fibre two-sphere, in the rest frame of \(y\), it is \(\mathrm{diag}\bigl(m^2,\tfrac13(2\ell^2-m^2),\tfrac13(2\ell^2-m^2),\tfrac13(2\ell^2-m^2)\bigr)\) (exact; checked on the explicit realisation). So the averaged kinetic tensor, which acts on fibre-constant functions, is positive definite for \(0<m^2<2\ell^2\), degenerate at \(m=0\) and \(m^2=2\ell^2\), and Lorentzian with the overall sign reversed for \(m^2>2\ell^2\). This is case (iii). It is the mechanism of~\cite{Steinacker2018BB}, where, for five matrices with the target form \(\mathrm{diag}(-\kappa^2,1,1,1,1)\), the averaged metric is positive definite everywhere for \(\kappa^2=1\), where one mass parameter suffices, and Lorentzian after a change of signature only for \(\kappa^2>1\), which needs two different mass parameters; the brane is Euclidean before the Big Bang there. In the Euclidean model,~\cite{LiaoMaeta2025} finds a solution formed by the six generators of \(\mathfrak{so}(1,3)\) whose effective metric is Euclidean, as (i) requires.

\emph{Coupling.} The map \(y^A=M^{A5}\), \(A=0,\dots,4\), onto the hyperboloid \(y_Ay^A=-\ell^2\) is equivariant under the \(\mathrm{SO}(4,1)\) that fixes the index \(5\), so it carries the symplectic measure to the invariant measure of the hyperboloid, which in the coordinates \(y^\mu\) is a constant times \(\mathrm d^4y/|m|\) on each of the parts \(m>0\) and \(m<0\). With the dilaton \(\rho^2=\rho_M|\det\gamma^{\mu\nu}|^{1/2}\) and coupling \(g_{\mathrm{YM}}^2=1/\rho^2\) of~\cite{Steinacker2026Weak} (its Eqs.~(88) and~(89)), the coupling is proportional to \(|m|^{-3}\) in the conjugate presentation, the dependence \(1/\sinh^3\tau\) of its Eq.~(89), now at every \(m\neq0\), and to \(|2\ell^2-m^2|^{-3/2}\) in the base presentation. In the conjugate presentation it diverges at \(m=0\), the only set where the kinetic tensor is not Lorentzian. There~\cite{Steinacker2026Weak} estimates the ratio of the two sides of the criterion as \(g_{\mathrm{YM}}/\sqrt2\) at late \(\tau\) (its Eq.~(100)), which places its semi-classical domain inside the Lorentzian one; that estimate does not extend to \(m\to0\). In the base presentation the coupling is finite on the whole positive definite domain, and no estimate of the fluctuations is given here.

These statements are for the Poisson limit, tree level, the background Feynman gauge for the commuting point, and the two four-matrix presentations of the orbit named, in the order \(N\to\infty\) first, as the Lorentzian member requires, then the \(\mathrm i\epsilon\) of the propagator, at fixed \(\gamma\) of the sign stated. The limits taken at finite \(N\) in the preceding paragraphs, with \(\epsilon\to0\) first, therefore do not apply to these backgrounds, and the conjugate presentation needs fixed \(\gamma>0\), which those paragraphs leave open. The criterion is not evaluated for the transverse fluctuations, the factor \(K\), the curved members of the first-order class or one loop. Nothing here selects the real form of the background algebra, the four-plane \(W\) or the background itself, and the non-degeneracy of the frame on \(m\neq0\) is a property of the supplied class, so Assumptions~\ref{ass:lorentzian}, \ref{ass:frame-lock} and~\ref{ass:nondeg} are not discharged; the kinetic tensor is not identified with the emergent metric of Chapter~\ref{ch:vierbein-emergence}. Problems~D and~E remain open and PROGRAMME-CONDITIONAL.
